% !TeX root = ../../Book3.tex
\section{Weierstrass 乘积定理}
\label{b3:sec:2003}

离散性是指定零点的唯一存在性障碍。即使没有任何有限幂和收敛，逐项增加初等因子的阶数仍可使尾部收敛。

\subsection{指定零点}
\label{b3:subsec:200301}

\begin{theorem}[Weierstrass：整平面版本]\label{b3:thm:20-weierstrass}
任意无有限聚点的复平面点列与指定正整数重数，可以实现为一个非恒零整函数的全部零点及重数。
\end{theorem}
\begin{proof}
把非零点按重数重复记为 \(a_n\)，则 \(|a_n|\to\infty\)。取
\(P(z)=z^m\prod_{n\ge1}E_n(z/a_n)\)。对固定 \(|z|\le R\)，充分大 \(n\) 时 \(|z/a_n|\le1/2\)，故对数尾项不超过 \(2\cdot2^{-n-1}\)，一致可和。正规乘积命题给整函数；有限头部与非零尾部分离，保证零点及重数正好指定。
\end{proof}
所得函数可能增长非常快，定理没有给阶或型的最优控制。

\subsection{典范乘积}
\label{b3:subsec:200302}

若零点的收敛指数有限，可取固定 \(p\) 的典范乘积，而不用变动的 \(E_n\)。这种选择减少指数修正并保留增长信息。

给定同一零点列，所有实现函数都形如 \(P e^g\)，其中 \(g\) 整。因为商在每个零点处可去且无零点，商的对数导数有整平面原函数，指数还原该商。加法常数由一个点的商值决定。

\subsection{零点序列}
\label{b3:subsec:200303}

一个点列若在有限点聚集，任何包含它的整函数都由恒等定理恒零；这说明“无有限聚点”既必要又充分。有限列给多项式，无穷列给真正的超越整函数。

指定重数并不指定导数系数。若还要在离散点给出函数值或有限阶导数，需要插值；下一节用主要部分给出这种额外自由度的构造。


若指定的零点极为缓慢地趋于无穷，例如 \(a_n=\log(n+1)\)，任何固定 \(p\) 的幂和都发散。Weierstrass 定理仍适用：对紧集 \(|z|\le R\)，从足够大的 \(n\) 起有 \(|z/a_n|\le1/2\)，变动阶数 \(p_n=n\) 使误差指数衰减。这个例子区分了“离散性保证存在”与“有限收敛指数保证有限属”。对应零点计数约为 \(e^r\)，所以任何实现它们的非恒零整函数都必须是无穷阶；这由上一章的零点密度限制直接得到。

\ProseParagraphBreak

相反，零点 \(a_n=n^2\) 的属零乘积具有较慢增长。考虑直接由幂级数定义的整函数
\[
F(z)=\sum_{k\ge0}\frac{(-1)^k\pi^{2k}z^k}{(2k+1)!}
=\frac{\sin(\pi\sqrt z)}{\pi\sqrt z}
\]
的后一表达式按前面级数解释，不需要在全平面选择平方根。更换平方根的符号不改变商，所以局部表达也相容。它的零点恰为正整数平方，每个为单零点；在负实轴 \(z=-r\) 上由 \(\sinh(\pi\sqrt r)\) 得增长下界，复平面上由指数表示得同阶上界，故阶为二分之一。一个非整数阶整函数因而不必靠抽象存在定理展示，可以直接从偶幂级数得到。

\ProseParagraphBreak

对于有限零点集，分解 \(f=P e^g\) 中 \(P\) 是普通多项式。若 \(f\) 又有多项式增长，\(g\) 的自由度实际上只能给常数因子；若只要求有限阶，则 \(g\) 可为多项式。若没有任何增长条件，\(g=e^z\) 等无穷阶例子仍可出现。零点集合本身无法区分这些函数，必须把无穷远条件作为独立数据。


\ProseParagraphBreak

构造的局部性可以用圆盘上的有限头部说明。给定 \(R\)，离散性使 \(|a_n|\le2R\) 的零点只有有限多个。它们的因子可能很大，但全是已固定的整函数；收敛只由其余尾部负责。对 \(|a_n|>2R\)，取 \(p_n=n\) 后
\[
\sup_{|z|\le R}|\log E_n(z/a_n)|\le2(R/|a_n|)^{n+1}\le2^{-n}.
\]
这给出了无需根据每个紧集重新排序的选择。所有半径使用同一个乘积，只是各自的有限头部长度不同。零点集合若在平面内有聚点，上述有限性马上失去；恒等定理同时说明非恒零整函数也不可能实现这类数据。

\ProseParagraphBreak

存在性没有给出规范化的唯一性。指定 \(P(0)=1\) 只消除了常数乘法自由度，仍可乘上 \(e^{az}\)；再指定 \(P'(0)=0\) 也仍可乘上 \(e^{bz^2}\)。只有增长约束或更多结构条件才会继续减少自由度。在有限阶一的情形，两个同零点函数分别有指数一次多项式的分解，因此它们的商为 \(e^{az+b}\)；这时指定原点值及一次对数导数才足够。若原点是零点，应先除去共同的原点零因子，再施加这些条件。
